% ============================================================================
% Fichier  : partie3/P03_C12_audit.tex
% Titre    : L'Audit selon l'ISO 19011
% Auteur   : MINKA MI NGUIDJOÏ Thierry Emmanuel
% Version  : 1.0 -- novembre 2026
% ============================================================================
% Sources : ISO 19011:2018 ; ISO/IEC 27007:2020 ; loi n° 2010/012 (art. 13 et
% 32, cités au chapitre 5). Programme d'audit de LiZbiZ : cas fictif.
% ============================================================================

\chapter{L'Audit selon l'ISO 19011}
\label{chap:audit}
\chaptermark{L'audit selon l'ISO 19011}

\epigraph{%
    « Auditer, ce n'est pas chercher des fautes~;\\
    c'est établir, preuves à l'appui, ce qui fonctionne et ce qui manque. »%
}{Principe du cours}

\niveauChapitre{Intermédiaire}{Chapitres~\ref{chap:gouvernance},
\ref{chap:droit} et~\ref{chap:traiter}~; cours INF3226.}

\begin{boxobjectifs}
\begin{itemize}
    \item Distinguer les audits de première, deuxième et troisième partie.
    \item Énoncer les sept principes de l'audit de l'ISO~19011.
    \item Établir un programme d'audit fondé sur les risques.
    \item Définir les rôles d'une équipe d'audit et les compétences
          attendues.
    \item Situer l'ISO/IEC~27007 pour l'audit d'un système de management de
          la sécurité.
\end{itemize}
\end{boxobjectifs}

\bigskip

% ============================================================================
\section*{Question directrice : que fallait-il auditer, et quand, chez LiZbiZ~?}
% ============================================================================

Le dernier audit de sécurité de LiZbiZ datait de 2023 et ne couvrait pas
le prestataire. Au printemps 2026, la mission d'AUDIX CONSEIL a examiné
cinq domaines du système d'information (cours INF3226). En septembre, le
comité d'audit demande à Mme~NKENG un \emph{programme} d'audit sur deux
ans, qui couvre en priorité ce qui présente le plus de risques, y compris
le dispositif de gestion des risques lui-même.

\medskip
\textbf{Plan du chapitre.}
\begin{enumerate}
    \item Les types d'audit (section~\ref{sec:c12-types}).
    \item Les principes de l'audit (section~\ref{sec:c12-principes}).
    \item Le programme d'audit (section~\ref{sec:c12-programme}).
    \item Les rôles et les compétences (section~\ref{sec:c12-roles}).
    \item L'audit d'un système de management de la sécurité (section~\ref{sec:c12-27007}).
\end{enumerate}

\separateur

% ============================================================================
\section{Les Types d'Audit}
\label{sec:c12-types}
% ============================================================================

\begin{boxdef}
\termecle{Audit} (ISO~19011)~: processus méthodique, indépendant et
documenté, permettant d'obtenir des preuves objectives et de les évaluer de
manière objective pour déterminer dans quelle mesure les critères d'audit
sont satisfaits~\cite{iso19011_2018}.
\end{boxdef}

\begin{tableaucours}{L{3.2cm}L{4.4cm}L{4.6cm}}{%
    \tableauHeader{Type} &
    \tableauHeader{Qui audite} &
    \tableauHeader{Exemple}
}
Première partie & L'organisation elle-même (audit interne) & Mme~NKENG
audite le dispositif de risques \\
Deuxième partie & Un client ou un donneur d'ordre & LiZbiZ audite SAHEL
INFOGÉRANCE, selon la clause du contrat \\
Troisième partie & Un organisme indépendant & Certification ISO/IEC~27001~;
audit de l'Agence au titre de l'article~32 \\
\end{tableaucours}
\vspace{-0.8em}
\begin{center}{\small\itshape Tableau~12.1. Les trois types d'audit}\end{center}

Les \termecle{critères d'audit} sont le référentiel auquel on compare les
preuves~: une norme, une politique interne, une clause de contrat, un
article de loi. Sans critère, il n'y a pas d'écart, seulement des opinions.

% ============================================================================
\section{Les Principes de l'Audit}
\label{sec:c12-principes}
% ============================================================================

\begin{tableaucours}{L{4.2cm}L{8.0cm}}{%
    \tableauHeader{Principe} &
    \tableauHeader{Ce qu'il exige}
}
Intégrité & Honnêteté, diligence, respect de la loi, impartialité \\
Présentation impartiale & Rapporter fidèlement, y compris les désaccords
non résolus \\
Conscience professionnelle & Diligence et jugement proportionnés à
l'importance de la tâche \\
Confidentialité & Protéger les informations obtenues \\
Indépendance & Ne pas auditer son propre travail~; rester libre de tout
conflit d'intérêts \\
Approche fondée sur la preuve & Des conclusions vérifiables, fondées sur un
échantillon d'informations disponibles \\
Approche fondée sur le risque & Concentrer l'audit sur ce qui compte pour
l'organisation et pour l'atteinte des objectifs de l'audit \\
\end{tableaucours}
\vspace{-0.8em}
\begin{center}{\small\itshape Tableau~12.2. Les sept principes de l'audit (ISO 19011:2018)}\end{center}

\begin{boxalert}
\textbf{Violation typique.} Confier l'audit du dispositif de risques à la
personne qui l'a conçu.

\smallskip
\textbf{Règle~:} Mme~MOUAFO ne peut pas auditer le registre qu'elle tient~;
Mme~NKENG le peut, parce qu'elle n'y a pas pris part
(Chapitre~\ref{chap:gouvernance}). Si l'auditeur a conseillé la
conception, il le déclare et l'on confie le sujet à un autre.
\end{boxalert}

% ============================================================================
\section{Le Programme d'Audit}
\label{sec:c12-programme}
% ============================================================================

\begin{boxdef}
\termecle{Programme d'audit}~: ensemble d'un ou plusieurs audits planifiés
pour une durée donnée et dirigés vers un but spécifique. Il précise les
objectifs, les risques et opportunités du programme, l'étendue, les
ressources, le calendrier et la manière de le surveiller.
\end{boxdef}

Un programme fondé sur les risques consacre plus d'audits, et plus tôt,
aux processus les plus exposés, en s'appuyant sur le registre des risques
(Chapitre~\ref{chap:matrices}).

\begin{tableaucours}{L{4.4cm}L{2.0cm}L{2.6cm}L{3.0cm}}{%
    \tableauHeader{Audit} &
    \tableauHeader{Risques} &
    \tableauHeader{Période} &
    \tableauHeader{Type}
}
Accès distants et prestataire & R-02 & Octobre 2026 & 1\textsuperscript{re}
et 2\textsuperscript{e} partie \\
Dispositif de gestion des risques & Tous & Octobre 2026 & 1\textsuperscript{re}
partie \\
Sauvegardes et reprise & R-01 & Janvier 2027 & 1\textsuperscript{re} partie \\
Données personnelles & R-07 & Mars 2027 & 1\textsuperscript{re} partie,
avec conseil \\
Continuité des agences & R-03, R-11 & Juin 2027 & 1\textsuperscript{re}
partie \\
Préparation à la certification ISO/IEC~27001 & SMSI & Septembre 2027 &
Audit à blanc \\
\end{tableaucours}
\vspace{-0.8em}
\begin{center}{\small\itshape Tableau~12.3. Programme d'audit 2026-2027 de LiZbiZ}\end{center}

\begin{boxscenario}{Opération LIZBIZ. Le programme validé}
\textbf{Constatation.} Le comité d'audit valide le programme le
30~septembre. Le premier audit, « accès distants et prestataire », combine
un audit interne et l'audit de SAHEL INFOGÉRANCE prévu par la nouvelle
clause du contrat.

\medskip
\textbf{Analyse} (confiance élevée). Le programme suit l'ordre des risques
du registre~: c'est l'approche fondée sur le risque. Il prévoit en outre
d'auditer le dispositif de risques lui-même, sans quoi un registre faux
orienterait faussement tout le programme.
\end{boxscenario}

% ============================================================================
\section{Les Rôles et les Compétences}
\label{sec:c12-roles}
% ============================================================================

\begin{tableaucours}{L{4.0cm}L{8.2cm}}{%
    \tableauHeader{Rôle} &
    \tableauHeader{Responsabilité}
}
Responsable du programme & Établit, met en œuvre, surveille et améliore le
programme \\
Responsable de l'équipe d'audit & Prépare et conduit un audit, répartit les
tâches, rédige le rapport \\
Auditeur & Recueille et évalue les preuves, rédige les constats \\
Expert technique & Apporte une compétence particulière, sous la
responsabilité d'un auditeur \\
Guide et observateur & Accompagnent l'équipe sans participer à
l'évaluation \\
\end{tableaucours}
\vspace{-0.8em}
\begin{center}{\small\itshape Tableau~12.4. Les rôles d'un audit}\end{center}

L'ISO~19011 décrit les compétences de l'auditeur~: comportement
professionnel (éthique, ouverture, observation, ténacité), connaissances
génériques (principes, méthodes, systèmes de management) et connaissances
propres à la discipline (ici, les risques et la sécurité de l'information).

% ============================================================================
\section{L'Audit d'un Système de Management de la Sécurité}
\label{sec:c12-27007}
% ============================================================================

L'\norme{ISO/IEC~27007:2020} complète l'ISO~19011 pour l'audit d'un système
de management de la sécurité de l'information~\cite{iso27007_2020}. Elle
précise ce qu'il faut examiner aux clauses~4 à~10 de l'ISO/IEC~27001
(contexte, leadership, planification dont l'appréciation des risques,
support, fonctionnement, évaluation, amélioration) et comment relier les
mesures de l'annexe~A aux risques.

\begin{boxobjet}{R-02 · Point d'audit}
\textbf{Scénario audité~:} usurpation de l'accès distant du prestataire.
\textbf{Critères~:} politique d'accès distant, clause du contrat
(octobre 2026), article~29, alinéa~2 de la loi n°~2010/012.
\textbf{Preuves attendues~:} configuration de la passerelle, journaux de
connexion, liste des comptes du prestataire, procès-verbal de l'audit du
prestataire. \textbf{Propriétaire~:} M.~NGUELE.
\end{boxobjet}

\begin{boxmetier}{Auditeur des systèmes d'information}
\textbf{Écrire les critères avant d'arriver.} L'auditeur liste, avant la
réunion d'ouverture, chaque critère et la preuve qui permettra de dire s'il
est satisfait. Ce qu'il ne peut pas relier à un critère n'ira pas dans le
rapport comme écart.
\end{boxmetier}

\separateur

\begin{boxresume}
\textbf{Ce qu'il faut retenir de ce chapitre~:}
\begin{itemize}
    \item \textbf{L'audit} compare des preuves objectives à des critères~;
          il est interne, de deuxième ou de troisième partie.
    \item \textbf{Sept principes}, dont l'indépendance, l'approche fondée
          sur la preuve et l'approche fondée sur le risque.
    \item \textbf{Le programme} suit le registre des risques et audite aussi
          le dispositif de risques.
    \item \textbf{Les rôles} sont distincts~; l'expert technique travaille
          sous la responsabilité d'un auditeur.
    \item \textbf{L'ISO/IEC~27007} adapte l'ISO~19011 au système de
          management de la sécurité.
\end{itemize}
\end{boxresume}

\begin{boxinfo}
\textbf{Points de contrôle de l'élève.} Je sais~: définir l'audit et les
critères d'audit~; distinguer les trois types~; citer les sept principes~;
construire un programme fondé sur les risques~; situer l'ISO/IEC~27007.
\end{boxinfo}

\bigskip

\begin{boxexo}[title={\ding{45}\quad Exercice~12.1. Quel type~? \niveaubase}]
Classez~: l'audit de LiZbiZ par son assureur~; l'audit de la paie par
Mme~NKENG~; l'audit de certification par un organisme accrédité.
\end{boxexo}

\begin{boxexo}[title={\ding{45}\quad Exercice~12.2. Les critères \niveaubase}]
Donnez trois critères d'audit possibles pour l'audit « sauvegardes et
reprise ».
\end{boxexo}

\begin{boxexo}[title={\ding{45}\quad Exercice~12.3. Principes \niveaubase}]
Quel principe un auditeur viole-t-il en omettant du rapport un désaccord
avec l'audité~?
\end{boxexo}

\begin{boxexo}[title={\ding{45}\quad Exercice~12.4. Indépendance \niveaumoyen}]
M.~KOUNOU propose d'auditer lui-même les sauvegardes, qu'il connaît
mieux que quiconque. Que répondez-vous~? Comment utiliser sa compétence~?
\end{boxexo}

\begin{boxexo}[title={\ding{45}\quad Exercice~12.5. Fondé sur le risque \niveaumoyen}]
Justifiez l'ordre du programme du Tableau~12.3 par le registre des risques
du Chapitre~\ref{chap:matrices}.
\end{boxexo}

\begin{boxexo}[title={\ding{45}\quad Exercice~12.6. Le prestataire \niveaumoyen}]
Quelles conditions doit réunir l'audit de SAHEL INFOGÉRANCE par LiZbiZ
pour être possible et utile~?
\end{boxexo}

\begin{boxexo}[title={\ding{45}\quad Exercice~12.7. L'expert \niveaumoyen}]
AUDIX CONSEIL fournit un expert des passerelles d'accès distant. Quel est
son rôle, et qui signe les constats~?
\end{boxexo}

\begin{boxexo}[title={\ding{45}\quad Exercice~12.8. Les risques du programme \niveaumoyen}]
Citez trois risques qui menacent le programme d'audit lui-même, et une
parade pour chacun.
\end{boxexo}

\begin{boxexo}[title={\ding{45}\quad Exercice~12.9. Un programme bancaire \niveauexpert}]
Construisez le programme d'audit sur deux ans de la BANQUE SANAGA,
fondé sur ses risques.
\end{boxexo}

\begin{boxexo}[title={\ding{45}\quad Exercice~12.10. Clauses 4 à 10 \niveauexpert}]
Pour chaque clause de 4 à 10 de l'ISO/IEC~27001, donnez une preuve que
l'auditeur demandera à LiZbiZ.
\end{boxexo}

\bigskip

\begin{boxinfo}
\textbf{Pour aller plus loin.}
\begin{itemize}
    \item ISO~19011~\cite{iso19011_2018}~; ISO/IEC~27007~\cite{iso27007_2020}.
\end{itemize}
\end{boxinfo}

\begin{boxinfo}
\textbf{Transition.} Le programme posé, il faut conduire chaque audit. Le
Chapitre~\ref{chap:conduire} traite des preuves, des échantillons et des
constats.
\end{boxinfo}
